\chapter{Exponential Families and the Operator Layer}
\label{ch:expfamily}

The belief and model law fibers of the main manuscript are collections of
normalized laws (\Cref{ch:geometry}, \eqref{eq:geo-law-fibers}), acted on by
pushforward (\eqref{eq:geo-pushforward-actions}); coordinates, densities, and
convex parameter sets are further choices. This chapter asks what a
finite-dimensional exponential family adds to a law fiber and what happens
when an algebraic map of its natural parameters leaves the natural domain. The
answer is that a natural-domain boundary point carries no law
(\Cref{prop:exp-domain-boundary-no-law}) and that a fixed point of a map on an
operator cone need not be a law (\Cref{prop:exp-fixed-point-no-law}), so that
pinning, quotienting, and retaining a mass coordinate are separate
constructions (\Cref{sec:exp-repairs}); inside the domain, relative entropy is
a Bregman divergence and affine mean constraints have a unique information
projection with a Pythagorean identity
(\Cref{prop:exp-affine-mean-projection}). The Gaussian realization of
\Cref{ch:generative-gaussian} is the case used throughout this supplement, the
Fisher charts of \Cref{ch:infogeometry} are its dual coordinates, and the
scale-free boundary operator of \Cref{ch:gaussian-renormalization} is an
instance of \Cref{prop:exp-fixed-point-no-law}.

This chapter derives no coarse-graining statement, and no Fisher or evidence
reading of one, from exponential-family structure alone: an operator arrow
needs a preserved parameter set, scale comparisons, and a declared rescaling,
while a law-level arrow needs normalized laws, compatible reference measures,
and finite normalizers at every scale. The boundary results below do not
exclude a quotient, a pin, a new reference measure, or a boundary stratum;
each is a different object that must be constructed. \status{NOT-CLAIMED}
Two questions remain open: the scaling of a retained mass coordinate
(\Cref{open:exp-mass-sector}) and the commutation of the information
projection with gauge actions and coarse parameter maps
(\Cref{open:exp-projection-commutation}).

\section{Regular finite-dimensional exponential families}
\label{sec:exp-family}

\definitionheading{Canonical family}{def:exp-canonical-family}
Let \(\nu\) be a nonzero sigma-finite measure on a standard Borel space
\((\mathsf Z,\mathscr Z)\), and let \(\mathsf T:\mathsf Z\to\mathsf V\) be a
measurable statistic into a finite-dimensional inner-product space
\(\mathsf V\). The log partition and the natural domain are
\begin{equation}
\mathsf A(\lambda)
=\log\int_{\mathsf Z}e^{\langle\lambda,\mathsf T(z)\rangle}\nu(dz),
\qquad
\mathsf N=\{\lambda\in\mathsf V:\mathsf A(\lambda)<\infty\},
\label{eq:exp-logpartition}
\end{equation}
and for \(\lambda\in\mathsf N\) the family member is
\begin{equation}
q_\lambda(z)
=\exp\{\langle\lambda,\mathsf T(z)\rangle-\mathsf A(\lambda)\},
\qquad
Q_\lambda(dz)=q_\lambda(z)\nu(dz).
\label{eq:exp-density}
\end{equation}
The vector \(\lambda\) is the natural parameter. No law is attached to a
parameter outside \(\mathsf N\).

\hypothesisheading{Regularity and minimality}{hyp:exp-regular-minimal}
Where a result names this hypothesis, the family is regular, so that
\(\mathsf N\) is open, and minimal, so that no nonzero \(a\in\mathsf V\) makes
\(\langle a,\mathsf T\rangle\) constant \(\nu\)-almost everywhere. Regularity
excludes natural-domain boundary points from the parameter set; minimality
excludes affine redundancy of the statistic.

\propositionheading{Log partition, cumulants, and Fisher information}{prop:exp-cumulants-fisher}
Assume \(\mathsf N\neq\varnothing\). (i) The extended-real function
\(\mathsf A\) is convex and lower semicontinuous on \(\mathsf V\), and
\(\mathsf N\) is convex. (ii) On \(\operatorname{int}\mathsf N\) the function
\(\mathsf A\) is infinitely differentiable, its derivatives are computed by
differentiating under the integral sign, and \(\mathsf T\) has moments of
every order, with
\begin{equation}
\nabla\mathsf A(\lambda)=\E_\lambda[\mathsf T],
\qquad
\nabla^2\mathsf A(\lambda)=\operatorname{Cov}_\lambda(\mathsf T).
\label{eq:exp-cumulants}
\end{equation}
Since the score of the family is \(\mathsf T-\nabla\mathsf A(\lambda)\), the
Hessian \(\nabla^2\mathsf A(\lambda)\) is the Fisher information in the
natural chart. (iii) If the family is minimal, then \(\mathsf A\) is strictly
convex on \(\mathsf N\), \(\nabla^2\mathsf A(\lambda)\) is positive definite at
every interior parameter, and \(\lambda\mapsto Q_\lambda\) is injective on
\(\mathsf N\) \citep[Theorems~7.1, 7.2 and~8.1 and Section~8.1]{BarndorffNielsen1978}.
\status{STANDARD}

\noindent\emph{Proof.}
The cited results concern a sigma-finite measure on a Euclidean space; they
apply to the image measure \(\mathsf T_\#\nu\) on \(\mathsf V\), because
\(\mathsf A\) is the logarithm of its Laplace transform and, for any
\(\lambda_0\in\mathsf N\), it has the positive density
\(e^{-\langle\lambda_0,t\rangle}\) with respect to the finite measure
\(e^{\langle\lambda_0,t\rangle}\mathsf T_\#\nu(dt)\). Minimality is the
statement that \(\mathsf T_\#\nu\) is not concentrated on an affine
hyperplane, and positive definiteness in (iii) then follows from
\(a^\top\nabla^2\mathsf A(\lambda)a=\operatorname{Var}_\lambda\langle a,\mathsf T\rangle\).
The domination behind (ii) is the elementary bound
\begin{equation}
|x|^r\leq \mathrm c_{r,\varepsilon}\left(e^{\varepsilon x}+e^{-\varepsilon x}\right),
\qquad r\in\{0,1,2\},\ x\in\R,
\label{eq:exp-domination}
\end{equation}
applied to each coordinate of \(\mathsf T\) at the vertices of a cube of side
\(4\varepsilon\) inside \(\mathsf N\) around an interior parameter.
\(\square\)

\propositionheading{Finite-domain interface to evidence-bound coordinates}{prop:exp-finite-domain-elbo-interface}
The map \((\lambda,z)\mapsto q_\lambda(z)\) is jointly measurable on
\(\mathsf N\times\mathsf Z\), and every law, density ratio, Fisher tensor, or
evidence-bound coordinate built from the canonical representation is defined
only for parameters in \(\mathsf N\). If two canonical representations share
one sigma-finite base measure \(\nu\), their density ratio is defined at every
pair of parameters taken each from its own natural domain. Consequently an
algebraic map that leaves \(\mathsf N\) ceases to define a probability law or
an evidence-bound coordinate, even when it remains in an operator cone.
\status{ESTABLISHED}

\noindent\emph{Proof.}
Lower semicontinuity of \(\mathsf A\)
(\Cref{prop:exp-cumulants-fisher}(i)) makes it Borel measurable, and
measurability of \(\mathsf T\) makes \((\lambda,z)\mapsto\langle\lambda,
\mathsf T(z)\rangle\) jointly measurable; exponentiation proves the first
claim. The normalizer is finite exactly on \(\mathsf N\) by
\eqref{eq:exp-logpartition}. With a common base measure, both densities in a
ratio are strictly positive \(\nu\)-densities of the form
\eqref{eq:exp-density}, so the ratio is defined at every pair of parameters in
the respective domains; representations with different statistics have
domains in different parameter spaces, so no intersection of domains is
formed. The last statement is a domain check against
\Cref{def:exp-partial-realization} and the extended evidence bound of
\Cref{def:elbo-extended}, which needs a normalized posterior slice of finite
positive mass. \(\square\)

An evidence-bound update additionally needs a positive finite evidence slice
at every compared parameter, and its classical split into two expected log
densities needs clauses (H3) and (H4) of \Cref{hyp:elbo-evidence-domain}.

\section{Operator and probability layers}
\label{sec:exp-layers}

\definitionheading{Partial realization}{def:exp-partial-realization}
A parametrized model consists of a parameter set \(\mathsf P\), a law domain
\(\mathsf P_{\mathrm{law}}\subseteq\mathsf P\), and a map
\begin{equation}
\mathfrak r:\mathsf P_{\mathrm{law}}\longrightarrow\mathcal B(\mathsf Z),
\qquad
\mathsf p\longmapsto Q_{\mathsf p},
\label{eq:exp-realization-map}
\end{equation}
into a collection \(\mathcal B(\mathsf Z)\) of normalized laws on
\(\mathsf Z\). A parameter outside \(\mathsf P_{\mathrm{law}}\) may still denote
an algebraic operator, but it denotes no law through \(\mathfrak r\).

\definitionheading{Operator layer}{def:exp-operator-layer}
Let \(\mathcal W\subseteq\mathsf V\) be a closed convex cone preserved by the
algebraic maps under study. Its probability layer is
\(\mathcal W\cap\mathsf N\). The cone is chosen for algebraic closure, not for
normalization, and a general law fiber carries no such cone.

\propositionheading{A natural-domain boundary carries no law}{prop:exp-domain-boundary-no-law}
Assume the family is regular and let \(\lambda^\star\in\partial\mathsf N\). Then
\begin{equation}
\int e^{\langle\lambda^\star,\mathsf T(z)\rangle}\nu(dz)=+\infty.
\label{eq:exp-boundary-diverges}
\end{equation}
For \(\lambda_0\in\mathsf N\), \(w=\lambda^\star-\lambda_0\), and
\(\lambda_t=(1-t)\lambda_0+t\lambda^\star\) with \(t\uparrow1\),
\begin{equation}
\mathsf A(\lambda_t)\to+\infty,\qquad
\langle w,\nabla\mathsf A(\lambda_t)\rangle\to+\infty,\qquad
\int_0^1\operatorname{Var}_{Q_{\lambda_t}}(\langle w,\mathsf T\rangle)dt
=+\infty.
\label{eq:exp-boundary-blowup}
\end{equation}
Hence the representation supplies at \(\lambda^\star\) no finite log
partition, mean parameter, Fisher metric, relative entropy from an endpoint
law, or evidence bound. \status{ESTABLISHED}

\noindent\emph{Proof.}
Regularity puts \(\lambda^\star\) outside the open set \(\mathsf N\), which is
\eqref{eq:exp-boundary-diverges}. Convexity of \(\mathsf N\)
(\Cref{prop:exp-cumulants-fisher}(i)) puts every \(\lambda_t\), \(t<1\), in
\(\mathsf N\), and lower semicontinuity forces the first limit. For
\(g(t)=\mathsf A(\lambda_t)\), convexity makes \(g'\) nondecreasing, and if
\(g'\) were bounded above then \(g\) would remain bounded, a contradiction.
The identities \(g'=\langle w,\nabla\mathsf A(\lambda_t)\rangle\) and
\(g''=\operatorname{Var}_{Q_{\lambda_t}}(\langle w,\mathsf T\rangle)\) of
\eqref{eq:exp-cumulants}, and integration of \(g''\), prove the remaining
limits. The last sentence restates that each listed object is built from
\(\mathsf A(\lambda^\star)\) or from \(Q_{\lambda^\star}\), neither of which
exists. \(\square\)

\propositionheading{An operator fixed point need not be a law}{prop:exp-fixed-point-no-law}
There is a regular canonical family with an operator cone \(\mathcal W\) and a
continuous map \(\psi:\mathcal W\to\mathcal W\) whose iterates from a point of
\(\mathcal W\cap\mathsf N\) all lie in \(\mathcal W\cap\mathsf N\), while its
unique fixed point lies on \(\partial\mathsf N\). One such family is
\(\mathsf Z=\R\), \(\mathsf T(z)=-z^2/2\), Lebesgue \(\nu\), for which
\begin{equation}
\mathsf A(\lambda)=-\tfrac12\log\lambda+\tfrac12\log2\pi,
\qquad
\mathsf N=(0,\infty),
\qquad
\mathcal W=[0,\infty),
\label{eq:exp-scale-witness}
\end{equation}
with \(\psi(\lambda)=\lambda/2\). \status{ESTABLISHED}

\noindent\emph{Proof.}
Starting at one gives \(\lambda_\ell=2^{-\ell}>0\), so every iterate is a law,
while \(\lambda_\ell\to0\), the unique fixed point, and
\(\mathsf A(\lambda_\ell)\to+\infty\). The constant integrand at
\(\lambda=0\) has infinite Lebesgue integral. \(\square\)

\section{Repairs of the boundary}
\label{sec:exp-repairs}

\definitionheading{Pinning}{def:exp-pinning}
Choose \(\lambda_{\mathrm m}\in\mathsf N\) and replace \(\mathcal W\) by
\(\lambda_{\mathrm m}+\mathcal W\); the pin does not flow.

\propositionheading{Pinning under a recession condition}{prop:exp-pinning-recession}
Write
\(\operatorname{rec}\mathsf N=\{\delta\in\mathsf V:\mathsf N+\delta\subseteq\mathsf N\}\)
for the recession cone of the convex, possibly nonclosed set \(\mathsf N\)
\citep[Theorem~8.1]{Rockafellar1970}. If
\(\lambda_{\mathrm m}\in\operatorname{int}\mathsf N\) and
\(\mathcal W\subseteq\operatorname{rec}\mathsf N\), then
\(\lambda_{\mathrm m}+\mathcal W\subseteq\operatorname{int}\mathsf N\).
\status{ESTABLISHED}

\noindent\emph{Proof.}
Let \(\delta\in\mathcal W\) and let \(\mathcal O\subseteq\mathsf N\) be an open ball
around \(\lambda_{\mathrm m}\). Since \(\delta\in\operatorname{rec}\mathsf N\),
the translated ball \(\mathcal O+\delta\) lies in \(\mathsf N\), so its center
\(\lambda_{\mathrm m}+\delta\) is interior. \(\square\)

A pin changes the family and adds a parameter that a map on the cone does not
move; removing it afterwards is a separate limit.

\definitionheading{Quotienting}{def:exp-quotienting}
Suppose \(\mathsf Z=\R^n\), \(\nu\) is Lebesgue measure, and the energy is
invariant along a nonzero linear subspace \(\mathcal V\). The quotient takes
the sample space \(\mathcal V^\perp\), its Lebesgue measure, and the restricted
statistic.

\propositionheading{Quotient normalization is not full-space normalization}{prop:exp-quotient-normalization}
Writing \(z=v+w\) with \(v\in\mathcal V\) and \(w\in\mathcal V^\perp\),
\begin{equation}
\int_{\R^n}e^{\langle\lambda,\mathsf T(z)\rangle}dz
=\operatorname{Leb}_{\mathcal V}(\mathcal V)
\int_{\mathcal V^\perp}e^{\langle\lambda,\mathsf T(w)\rangle}dw
=+\infty.
\label{eq:exp-quotient-split}
\end{equation}
The quotient law is normalized exactly when the second factor is finite. In
particular, a Moore--Penrose inverse and a pseudodeterminant inserted into the
full-space density formula compute the finite quotient expression, whose
dimension and reference measure differ, and leave the full-space integral
infinite. \status{ESTABLISHED}

\noindent\emph{Proof.}
The orthogonal change of variables has unit Jacobian and the integrand does
not depend on \(v\), so Tonelli's theorem gives
\eqref{eq:exp-quotient-split}, whose first factor is infinite. A
pseudoinverse and a pseudodeterminant evaluate the Gaussian integral on
\(\mathcal V^\perp\) only, which is the second factor. \(\square\)

\definitionheading{Retaining a mass sector}{def:exp-mass-sector}
Enlarge the state by a coercive parameter that remains in the natural domain
under the scale map. This is neither pinning nor quotienting: the mass is a
running coordinate, and its scaling law is part of the declared scale map.

\openproblemheading{Scaling of a retained mass sector}{open:exp-mass-sector}
Pinning keeps an orbit interior under the recession condition of
\Cref{prop:exp-pinning-recession}, and quotienting normalizes exactly when
the quotient integral of \eqref{eq:exp-quotient-split} is finite. For a
retained mass coordinate (\Cref{def:exp-mass-sector}) nothing here fixes its
law under a scale map, and so nothing decides whether its iterates stay in
\(\mathsf N\) or approach \(\partial\mathsf N\). The question is settled for a
given scale map by an explicit induced law on the mass coordinate together
with a proof that its orbit stays in \(\mathsf N\), or reaches the boundary,
under stated hypotheses on the map.

\section{Relative entropy and affine information projection}
\label{sec:exp-information-projection}

For a regular minimal family write
\(\boldsymbol\mu_\lambda=\nabla\mathsf A(\lambda)\) for the mean parameter and
\(\mathsf A^\ast\) for the convex conjugate of \(\mathsf A\). For a convex
function \(\phi\) differentiable at \(y\), let
\(D_\phi(x,y)=\phi(x)-\phi(y)-\langle\nabla\phi(y),x-y\rangle\) be its Bregman
divergence \citep{Bregman1967}.

\propositionheading{Bregman form of relative entropy}{prop:exp-kl-bregman}
Under \Cref{hyp:exp-regular-minimal}, for \(\lambda,\lambda'\in\mathsf N\),
\begin{equation}
\KL(Q_{\lambda'}\Vert Q_\lambda)
=\mathsf A(\lambda)-\mathsf A(\lambda')
-\langle\nabla\mathsf A(\lambda'),\lambda-\lambda'\rangle
=D_{\mathsf A}(\lambda,\lambda'),
\label{eq:exp-kl-bregman}
\end{equation}
and
\begin{equation}
\KL(Q_{\lambda'}\Vert Q_\lambda)
=D_{\mathsf A^\ast}(\boldsymbol\mu_{\lambda'},\boldsymbol\mu_\lambda).
\label{eq:exp-kl-dual-bregman}
\end{equation}
This is the exact divergence of the family; it is not in general a
finite-displacement Fisher quadratic. \status{ESTABLISHED}

\noindent\emph{Proof.}
The log density ratio is
\(\log(q_{\lambda'}/q_\lambda)=\langle\lambda'-\lambda,\mathsf T\rangle
-\mathsf A(\lambda')+\mathsf A(\lambda)\), and its expectation under
\(Q_{\lambda'}\), with \(\E_{\lambda'}\mathsf T=\nabla\mathsf A(\lambda')\)
from \eqref{eq:exp-cumulants}, is \eqref{eq:exp-kl-bregman}; this is
\citet[Proposition~6.3]{Brown1986}. The function \(\mathsf A\) is closed by
\Cref{prop:exp-cumulants-fisher}(i), has the open domain \(\mathsf N\), is
differentiable there, and is strictly convex by
\Cref{prop:exp-cumulants-fisher}(iii). The subgradients of \(\mathsf A^\ast\)
at \(\boldsymbol\mu_\lambda\) are therefore exactly the solutions of
\(\nabla\mathsf A(\lambda'')=\boldsymbol\mu_\lambda\), namely \(\lambda\)
alone, so \(\mathsf A^\ast\) is differentiable there with
\(\nabla\mathsf A^\ast(\boldsymbol\mu_\lambda)=\lambda\) and
\(\mathsf A^\ast(\boldsymbol\mu_\lambda)=\langle\lambda,\boldsymbol\mu_\lambda\rangle
-\mathsf A(\lambda)\) \citep[Theorems~23.5 and~25.1]{Rockafellar1970}.
Substituting these identities at \(\lambda\) and \(\lambda'\) into the
definition of \(D_{\mathsf A^\ast}\) gives \(D_{\mathsf A}(\lambda,\lambda')\),
which is \eqref{eq:exp-kl-dual-bregman}. \(\square\)

\propositionheading{Affine mean constraints attain a unique projection}{prop:exp-affine-mean-projection}
Assume the extended \(\mathsf A\) is Legendre, as in a regular minimal
canonical family, and let
\(\mathcal M=\operatorname{int}\operatorname{dom}\mathsf A^\ast\). If
\(\mathscr A\) is a closed affine subspace with
\(\mathscr A\cap\mathcal M\neq\varnothing\), then for every
\(\boldsymbol\mu_\lambda\in\mathcal M\) there is a unique
\(\widehat{\boldsymbol\mu}\in\mathscr A\cap\mathcal M\) minimizing
\begin{equation}
\boldsymbol\mu\longmapsto
D_{\mathsf A^\ast}(\boldsymbol\mu,\boldsymbol\mu_\lambda).
\label{eq:exp-mean-projection}
\end{equation}
Its natural parameter is
\(\widehat\lambda=\nabla\mathsf A^\ast(\widehat{\boldsymbol\mu})\), and every
natural parameter \(\lambda'\) with \(\boldsymbol\mu_{\lambda'}\in\mathscr A\cap\mathcal M\)
satisfies the Pythagorean identity
\begin{equation}
\KL(Q_{\lambda'}\Vert Q_\lambda)
=\KL(Q_{\lambda'}\Vert Q_{\widehat\lambda})
+\KL(Q_{\widehat\lambda}\Vert Q_\lambda),
\label{eq:exp-projection-pythagorean}
\end{equation}
so \(Q_{\widehat\lambda}\) is the unique minimizer of
\(\KL(\cdot\Vert Q_\lambda)\) over the members of the family whose mean
parameters lie in \(\mathscr A\). \status{ESTABLISHED}

\noindent\emph{Proof.}
The objective in \eqref{eq:exp-mean-projection} is the Bregman divergence of
the Legendre function \(\mathsf A^\ast\) from the interior point
\(\boldsymbol\mu_\lambda\). For such a function and a closed convex feasible
set meeting the interior of its domain, the minimizer exists, is unique, and
lies in that interior \citep[Theorem~3.12(iii)]{BauschkeBorwein1997}. Write
\(\mathscr A=\{\boldsymbol\mu:\mathsf D\boldsymbol\mu=\mathsf b\}\). The
objective is differentiable at the interior minimizer with gradient
\(\nabla\mathsf A^\ast(\widehat{\boldsymbol\mu})-\lambda\), which must
annihilate the direction space \(\ker\mathsf D\); conversely, by convexity, a
feasible interior point at which the gradient annihilates \(\ker\mathsf D\) is
a minimizer. Since \((\ker\mathsf D)^\perp=\operatorname{im}\mathsf D^\top\),
the solution is characterized by
\begin{equation}
\mathsf D\widehat{\boldsymbol\mu}=\mathsf b,\qquad
\nabla\mathsf A^\ast(\widehat{\boldsymbol\mu})-\lambda+\mathsf D^\top\omega=0
\label{eq:exp-projection-kkt}
\end{equation}
for some multiplier \(\omega\), which proves the natural-parameter formula;
\citet{Bauschke2003} gives the dual characterization of Bregman projections
onto affine subspaces. For the Pythagorean identity, write
\(\boldsymbol\mu'=\boldsymbol\mu_{\lambda'}\) and expand the three divergences
by \eqref{eq:exp-kl-bregman}; the values of \(\mathsf A\) cancel and
\[
\KL(Q_{\lambda'}\Vert Q_\lambda)
-\KL(Q_{\lambda'}\Vert Q_{\widehat\lambda})
-\KL(Q_{\widehat\lambda}\Vert Q_\lambda)
=\langle\widehat{\boldsymbol\mu}-\boldsymbol\mu',\lambda-\widehat\lambda\rangle.
\]
By \eqref{eq:exp-projection-kkt}, \(\lambda-\widehat\lambda=\mathsf D^\top\omega\),
while feasibility gives \(\mathsf D(\widehat{\boldsymbol\mu}-\boldsymbol\mu')=0\),
so the pairing vanishes. The middle divergence is nonnegative and vanishes
only when \(Q_{\lambda'}=Q_{\widehat\lambda}\), which gives the minimizing
property, and injectivity of \(\lambda\mapsto Q_\lambda\)
(\Cref{prop:exp-cumulants-fisher}(iii)) gives uniqueness. For a general
closed convex feasible set only the inequality form of
\eqref{eq:exp-projection-pythagorean} holds
\citep[Proposition~3.16]{BauschkeBorwein1997}. \(\square\)

\remarkheading{Feasibility, closedness, and convexity are needed}{rem:exp-projection-hypotheses}
Take the Bernoulli family with mean domain \((0,1)\) and target mean
\(\tfrac12\). The feasible set \(\{0\}\) misses the mean domain, so no
projection exists; on the open feasible interval \((0,\tfrac14)\) the
divergence is strictly decreasing and its infimum, attained only at the
excluded endpoint \(\tfrac14\), is finite and not attained; and the closed
nonconvex set \(\{\tfrac14,\tfrac34\}\) has two minimizers, because
\(\KL(\operatorname{Bern}(\mu)\Vert\operatorname{Bern}(\tfrac12))\) is
symmetric under \(\mu\mapsto1-\mu\). \status{ESTABLISHED}

\noindent\emph{Proof.}
The divergence \(\mu\log(2\mu)+(1-\mu)\log(2(1-\mu))\) has derivative
\(\log(\mu/(1-\mu))<0\) for \(\mu<\tfrac12\), which gives the interval claim,
and the symmetry is visible in the formula. \(\square\)

\openproblemheading{Commutation of the information projection}{open:exp-projection-commutation}
\Cref{prop:exp-affine-mean-projection} proves existence, uniqueness, and the
Pythagorean identity of the projection onto an affine mean constraint. It is
not proved whether the projection is equivariant under a gauge action on the
natural parameters, or whether it commutes with a coarse parameter map. The
question is settled by an intertwining theorem that states hypotheses on the
action or the map under which the two projections correspond, or by a
counterexample within a regular minimal family.
